\section*{September 8, 2026: Why the treated mean falls after 2016}
\textit{Drafted by Codex.}

Within the fixed quarter-mile sample, the treated mean falls from
\$338,890 in 2016 to \$319,898 in 2017 and \$293,567 in 2018, a 13.4 percent
2016--18 decline. The control mean rises from \$193,910 to \$229,863.
These are transaction prices in 2022 dollars; they do not follow fixed homes.

Changing neighborhood shares explains the direction of the treated decline.
Holding each of the 29 treated sites at its 2016 share of sales yields a 2018
mean of \$391,580, 15.5 percent above 2016. All 29 sites have observations in
each of 2016--18, so this exercise drops no treated sites. It holds geography
weights fixed, not property quality or the mix of homes sold within each site.
The actual sample and its transaction weights remain unchanged.

Peabody provides the clearest example: qualifying sales fall from 28 to eight,
while their average price rises from \$728,787 to \$1,011,699. Near North sales
fall from 17 to 11, Lafayette from 33 to 25, and King from 24 to 15.
Meanwhile West Pullman sales rise from seven to 31; its 2018 mean is only
\$50,528. Omitting Peabody alone reduces the aggregate dollar decline from
\$45,324 to \$27,494 but does not eliminate it. No single-site omission
eliminates the 2016--18 decline.

The upper tail changes mainly in its frequency. Sales above \$500,000 fall
from 103 of 369 transactions to 83 of 413. The mean among sales below that
fixed threshold is almost unchanged (\$151,276 versus \$151,345), and the
number above \$1 million is 24 versus 25. The 75th percentile falls from
\$530,216 to \$402,096; the 95th falls much less, from \$1,070,043 to
\$1,028,779. Broad property types alone do not settle the issue: both house
and apartment-building means fall, but their neighborhood composition can
change within each type.

The decline in Peabody activity precedes price cleaning: recorded transactions
in the same geography fall from 54 to 25; market-screened sales fall from 39
to 12, and the final sample from 28 to eight. Cleaning therefore reinforces
but does not originate this site's loss of weight. Both treatment groups
have transactions in all twelve months of 2016, 2017, and 2018. These checks
do not establish complete coverage of all actual market transactions.

Contemporaneous market evidence does not suggest a comparable citywide annual
price collapse. Illinois REALTORS reports that Chicago's 2018 sales fell
4.4 percent while its annual nominal median rose 1.8 percent. This covers
single-family homes and condos, a different universe from our sample.\footnote{\href{https://www.illinoisrealtors.org/blog/illinois-home-sales-lower-in-december-market-closes-out-a-largely-stable-2018/}{Illinois REALTORS, January 22, 2019.}}
A local event worth recording is Rowe's September 2017 expansion at the former
Peabody annex, following plans announced in 2015. The evidence does not link
that expansion to the drop in nearby sales.\footnote{\href{https://www.dnainfo.com/chicago/20170915/west-town/charter-school-once-packed-gills-shows-off-new-2nd-floor/}{DNAinfo, September 15, 2017, contemporaneous reporting recovered in indexed text.}}
The demonstrated explanation is changing sales composition; why activity
shifted across these particular neighborhoods remains unresolved.

\small Reproduction: \texttt{make ../output/late\_decline.pdf} in
\texttt{tasks/audits/home\_data\_overview/code}. The saved
\texttt{late\_decline.rds} matches all 11,600 packet transactions and annual
means. Source snapshots are unchanged from \texttt{e76f4dd}; scripts are the
September 8 working-tree extension. \normalsize
\clearpage
\includegraphics[page=1,width=\textwidth]{../input/late_decline.pdf}

\clearpage
\subsection*{Sales by pre-closure neighborhood price}
Each school site is classified using its median transaction price over
2008--2012 in 2022 dollars. The common cutoff is the median of those site
medians across both treatment groups, \$65,050. Sites at or below the cutoff
form the lower-price category. Classification stays fixed; the sale price of
an individual transaction never determines its category. The 77 sites with
sales all have pre-period observations, so all 11,600 transactions are retained.
The lower category contains 14 treated and 25 control sites; the higher
category contains 15 treated and 23 controls.

Between 2016 and 2018, treated sales rise from 116 to 182 in the lower category
and fall from 253 to 231 in the higher category. Control sales rise from 433
to 550 and from 349 to 403, respectively. The initially cheaper neighborhoods
therefore account for a growing share of treated sales. These are counts,
not turnover rates; the numbers of sites and housing stocks differ across groups.

\includegraphics[width=\textwidth]{../input/sales_by_pre_price.pdf}

\small Reproduction: \texttt{make ../output/sales\_by\_pre\_price.pdf} in
\texttt{tasks/audits/home\_data\_overview/code}. The classification and annual
counts are saved with the late-decline summaries and included in their standard
report. \normalsize
